\section{Trazabilidad de datos, cálculos y pruebas}
\label{app:trazabilidad}

Las rutas de este apéndice se expresan con respecto a la raíz del repositorio.
El punto de entrada de la evidencia del artículo es
\codepath{version_2/validation/paper07_chua/}. En esa carpeta,
\codepath{manifest.json} relaciona cada afirmación numérica con sus archivos
de origen y \codepath{evidence_manifest.json} comprueba que el paquete
compacto esté completo e íntegro. La carpeta
\codepath{evidence/} contiene los datos necesarios para auditar los resultados
archivados sin depender de una figura ni de una transcripción manual. Esta
integridad no implica que cada campaña histórica pueda reintegrarse desde una
copia limpia.

La trazabilidad sigue siempre la misma cadena:
\[
\text{parámetros y contrato}
\longrightarrow
\text{semilla}
\longrightarrow
\text{trayectoria}
\longrightarrow
\text{sondeo}
\longrightarrow
\text{tabla o figura}.
\]
Los archivos JSON fijan parámetros, tolerancias y decisiones; los CSV
conservan etapas o resultados por condición inicial; y los NPZ almacenan
bancos de parámetros, semillas o trayectorias cuando es necesario preservar
arreglos completos.

\subsection{Cadena del sistema no suave}

La información del caso de \cref{eq:chua_no_suave} se encuentra en
\codepath{version_2/validation/paper07_chua/evidence/nonsmooth_corrected/}.
La lectura recomendada es la siguiente:

\begin{enumerate}
  \item \codepath{candidate_and_reference.json} registra la raíz sesgada
  seleccionada, la semilla modal y la trayectoria de referencia;
  \item \codepath{numerical_contract.json} fija \(q\), el paso, los
  horizontes, el integrador ABM--PECE y la política de memoria;
  \item \codepath{continuation_stages.csv} conserva, en orden, los 21 valores
  de \(\eta\) y el estado alcanzado en cada etapa;
  \item \codepath{target_reproduction_runs.csv} documenta las reintegraciones
  usadas para comprobar la atracción local del conjunto objetivo;
  \item \codepath{result.json} reúne la semilla, el extremo de continuación y
  el diagnóstico postransitorio.
\end{enumerate}

El sondeo volumétrico usado en las tablas del reporte se conserva en
\codepath{nonsmooth_corrected/extended_first_contact_clean/}. Dentro de esa
carpeta, \codepath{extended_numerical_contract.json} fija el contrato,
\codepath{extended_probe_plan.csv} enumera los radios y tamaños de muestra,
\codepath{extended_probe_runs.csv} registra una fila por condición inicial y
\codepath{extended_probe_summary.csv} contiene los conteos agregados.
\codepath{extended_result.json} separa las \(7\,200\) pruebas del contrato
local \(r\leq0.01\) de las \(10\,200\) pruebas de la extensión macroscópica,
comprueba el total de \(17\,400\) y registra el corte después de completar
\(r=0.3\), primer radio con contactos.
\codepath{target_cloud_nn_sample.csv} fija la muestra del conjunto objetivo
empleada para medir contacto.

La integridad de esta evidencia fijada se comprueba mediante
\codepath{version_2/validation/paper07_chua/scripts/
sync_paper07_evidence.py}. El programa verifica el manifiesto compacto y, con
la opción correspondiente, su paridad con las fuentes canónicas. Esta
comprobación valida los artefactos conservados; no se presenta como una nueva
ejecución de la campaña numérica completa.

\subsection{Cadena del sistema con arctangente}

El registro canónico actualmente promovido para \texttt{c590} se encuentra en
\codepath{version_2/validation/chua_fractional_arctan_c590/}. El manifiesto
\codepath{version_2/validation/paper07_chua/manifest.json} conserva el origen de
la semilla como búsqueda entera acotada y refinamiento de Caputo, y el nombre
\texttt{c590} identifica el índice de aquella búsqueda local. El paquete
compacto valida los resultados retenidos, pero no contiene ya los programas
históricos de regeneración de las exploraciones global y local; por tanto, no
se afirma que esa etapa pueda repetirse desde este apéndice.

Los sondeos sobre superficies esféricas se encuentran en
\codepath{version_2/validation/chua_fractional_arctan_c590/}. Para las tablas
de este reporte se emplean, en orden,
\codepath{hiddenness_scaled_rows.csv},
\codepath{hiddenness_r003_rows.csv},
\codepath{hiddenness_r010_rows.csv},
\codepath{hiddenness_r030_rows.csv} y
\codepath{hiddenness_r100_rows.csv},
\codepath{hiddenness_r200_rows.csv}, junto con sus archivos
\codepath{*_run_config.json}. Esas seis corridas reúnen \(13\,500\)
condiciones iniciales: \(0/8\,400\) contactos hasta \(r=0.3\),
\(22/2\,400\) en \(r=1\) y \(588/2\,700\) en \(r=2\), con 131
divergencias en este último bloque.

Las trayectorias representativas fijadas y su manifiesto se conservan en
\codepath{version_2/validation/paper07_chua/evidence/probe_story_trajectories/};
son la fuente de los gráficos que ilustran los destinos de las sondas y forman
parte del manifiesto de integridad.

\subsection{Controles bibliográficos y validación algebraica}

Los controles de las referencias se mantienen separados de los dos resultados
propuestos:

\begin{itemize}
  \item Leonov--Kuznetsov:
  \codepath{version_2/validation/published_cases/kuznetsov2017_chua_integer.yaml}
  y \codepath{version_2/validation/reference_cases/chua_integer_q1/};
  \item Danca:
  \codepath{version_2/validation/published_cases/danca2017_chua_fractional_saturation.yaml};
  \item Wu et al.:
  \codepath{version_2/validation/published_cases/wu2023_chua_fractional_arctan.yaml}
  y
  \codepath{version_2/validation/reference_cases/fractional_chua_arctan_wu2023/}.
\end{itemize}

\codepath{version_2/validation/published_reference_coverage.json} indica qué
ecuaciones, parámetros, condiciones iniciales y procedimientos se verifican
en cada control. La comprobación simbólica se describe en
\cref{app:wolfram}; en particular, el archivo específico de
\texttt{c590} usa el valor \(\rho\neq1\) de \cref{eq:parametros_c590}.

\subsection{Pruebas automatizadas relevantes}

Las pruebas se agrupan por la propiedad científica que protegen:

\begin{table}[H]
\centering
\caption{Pruebas que enlazan la implementación con el procedimiento descrito.}
\label{tab:pruebas_trazabilidad}
\begin{tabularx}{\linewidth}{L{0.33\linewidth}Y}
\toprule
Prueba & Propiedad comprobada\\
\midrule
\codepath{version_2/tests/test_summarize_c590_hiddenness.py}
& Agregación de las seis corridas de \texttt{c590}, conteos por radio y frontera
de la afirmación finita.\\
\codepath{version_2/tests/test_continuation_memory_policy.py} y
\codepath{version_2/tests/test_fractional_memory_validation.py}
& Conservación de la memoria causal y consistencia del integrador
fraccionario.\\
\codepath{version_2/tests/test_published_case_reproduction.py},
\codepath{version_2/tests/test_published_continuation_comparison.py} y
\codepath{version_2/tests/test_published_validation_coverage.py}
& Separación y cobertura de los controles de Leonov--Kuznetsov, Danca y Wu.\\
\codepath{version_2/tests/test_official_wolfram_artifacts.py}
& Correspondencia entre los casos algebraicos, sus parámetros y sus salidas.\\
\codepath{version_2/tests/test_sync_paper07_evidence.py}
& Completitud del paquete compacto de evidencia del artículo.\\
\bottomrule
\end{tabularx}
\end{table}

Desde \codepath{version_2/}, una comprobación corta de procedencia e integridad
se realiza con:
\begin{verbatim}
python validation/paper07_chua/scripts/sync_paper07_evidence.py \
  --verify --verify-sources
python -m pytest -q tests/test_summarize_c590_hiddenness.py \
  tests/test_sync_paper07_evidence.py \
  tests/test_official_wolfram_artifacts.py
\end{verbatim}
Cada instrucción larga puede escribirse en una sola línea. Esta secuencia
comprueba la asociación entre parámetros, semillas, contratos y archivos
canónicos. No equivale a reejecutar las campañas integrales que produjeron los
artefactos fijados.
